\chapter{Výsledky}

V~tejto kapitole predkladáme podrobnú analýzu údajov získaných počas nášho výskumu. Zameriavame sa na porovnanie dvoch skupín pacientov s~reumatoidnou artritídou: experimentálnej skupiny (ITF + kinezioterapia) a~kontrolnej skupiny (samostatná kinezioterapia). Výsledky sme rozdelili na analýzu subjektívneho vnímania klinického stavu prostredníctvom dotazníka a~objektívne zhodnotenie mobility pomocou goniometrických meraní.

\section{Analýza výsledkov subjektívneho hodnotenia}

Na získanie údajov o~subjektívnom stave pacientov sme použili dotazník vlastnej konštrukcie, ktorý sledoval 10 parametrov (bolesť v~pokoji a~pri záťaži, ADL, ranná stuhnutosť, únava, jemná motorika, psychosociálne aspekty). Maximálne dosiahnuteľné skóre bolo 100 bodov, pričom vyššia hodnota indikovala horší klinický stav.

\begin{figure}[H]
    \centering
    % \includegraphics[width=0.8\textwidth]{obrazky/graf_dotaznik_celkove.png}
    \caption{Porovnanie celkového skóre dotazníka pred a~po terapii}
    \label{fig:graf_dotaznik}
\end{figure}

V~\textbf{experimentálnej skupine} (n=50), ktorej bola aplikovaná kombinovaná liečba (ITF + kinezioterapia), bola pred začatím terapie priemerná hodnota celkového skóre 59,2~$\pm$~18,56 bodu. Po absolvovaní mesiaca liečby kleslo priemerné skóre na 49,22~$\pm$~18,52 bodu. U~pacientov sme zaznamenali štatisticky významné zlepšenie vo všetkých sledovaných parametroch dotazníka, najvýraznejšie v~oblasti ranného stuhnutia a~bolesti pri chôdzi.

V~\textbf{kontrolnej skupine} (n=50), liečenej samostatnou kinezioterapiou, bola vstupná priemerná hodnota 50,35~$\pm$~18,50 bodu. Po mesiaci rehabilitácie sa priemerné skóre znížilo na 46,06~$\pm$~18,52 bodu. Hoci aj v~tejto skupine došlo k~pozitívnemu posunu, miera zlepšenia bola v~porovnaní s~experimentálnou skupinou nižšia. Pri porovnaní oboch skupín vykázala experimentálna skupina s~pridanou elektroterapiou výraznejší pokles v~celkovom bodovom hodnotení, čo svedčí o~vyššom analgetickom a~funkčnom benefite kombinovanej liečby.

\section{Analýza výsledkov goniometrických meraní}

Objektívne hodnotenie mobility sme realizovali meraním rozsahov pohybov v~postihnutých kĺboch hornej končatiny. Zamerali sme sa na kĺby, ktoré sú pri reumatoidnej artritíde najčastejšie limitované zápalovým procesom.

\subsection{Hodnotenie mobility ramenného kĺbu}

V~grafe \ref{fig:graf_rameno} sú zobrazené namerané hodnoty aktívneho rozsahu pohybu (flexie) v~ramennom kĺbe pred a~po absolvovaní terapeutického programu.

\begin{figure}[H]
    \centering
    % \includegraphics[width=0.8\textwidth]{obrazky/graf_gonio_rameno.png}
    \caption{Aktívny rozsah pohybu v~ramennom kĺbe pred a~po terapii}
    \label{fig:graf_rameno}
\end{figure}

V~\textbf{experimentálnej skupine} dosahoval pred terapiou priemerný rozsah flexie v~ramene 112,0~$\pm$~5,36 stupňa. Po aplikácii interferenčných prúdov a~cvičenia sa rozsah zvýšil na priemerných 122,26~$\pm$~5,42 stupňa, čo predstavuje priemerné zlepšenie o~viac ako 10 stupňov. V~\textbf{kontrolnej skupine} bol počiatočný priemerný rozsah 111,9~$\pm$~5,6 stupňa a~po mesiaci cvičenia sa zvýšil na 117,2~$\pm$~5,9 stupňa. Výsledky potvrdzujú, že pacienti liečení kombináciou ITF a~kinezioterapie dosiahli v~oblasti ramennej mobility štatisticky významnejší progres než pacienti s~izolovanou pohybovou liečbou.

\subsection{Hodnotenie mobility zápästného a~MCP kĺbov}

Zápästia a~drobné kĺby ruky sú kľúčové pre jemnú motoriku a~sebaobsluhu. Priemerné hodnoty mobiliy v~týchto segmentoch pred a~po liečbe uvádzame v~tabuľke \ref{tab:vysledky_ruka}.

\begin{table}[H]
    \centering
    \caption{Porovnanie rozsahov pohybu zápästia a~MCP kĺbov (stupne)}
    \label{tab:vysledky_ruka}
    \begin{tabular}{lcccc}
        \toprule
         & \multicolumn{2}{c}{\textbf{Exp. skupina (ITF+K)}} & \multicolumn{2}{c}{\textbf{Kontr. skupina (K)}} \\
        \textbf{Parameter} & \textbf{Pred} & \textbf{Po} & \textbf{Pred} & \textbf{Po} \\
        \midrule
        Zápästný kĺb (D/P flexia) & 41,46 ± 4,2 & 51,74 ± 4,2 & 45,48 ± 4,3 & 50,10 ± 4,4 \\
        MCP kĺby (flexia) & 66,78 ± 7,4 & 76,48 ± 7,1 & 70,02 ± 6,8 & 74,80 ± 6,5 \\
        MCP kĺby (extenzia) & 14,16 ± 3,2 & 20,40 ± 3,5 & 14,20 ± 3,1 & 18,40 ± 3,3 \\
        \bottomrule
    \end{tabular}
\end{table}

V~oblasti \textbf{MCP kĺbov} sme sa zamerali najmä na extenziu, ktorá býva pri RA často limitovaná v~dôsledku zápalových zmien väzivového aparátu. V~experimentálnej skupine sa priemerná extenzia zvýšila zo 14,16 stupňa na 20,40 stupňa, čo pacientom umožnilo lepší úchop a~funkciu ruky. U~kontrolnej skupiny bol nárast miernejší (na 18,40 stupňa). Zlepšenie v~zápästnom kĺbe bolo v~experimentálnej skupine taktiež výraznejšie (zlepšenie o~10,28 stupňa oproti 4,62 stupňa v~kontrolnej skupine), čo potvrdzuje aditívny relaxačný a~protiopuchový vplyv interferenčných prúdov.

\section{Verifikácia stanovených hypotéz}

Na základe analýzy a~štatistického spracovania výsledkov môžeme pristúpiť k~vyhodnoteniu stanovených výskumných hypotéz.

\textbf{H1: Potvrdená.} Predpoklad o~tom, že pacienti liečení kombinovanou terapiou uvádzajú po liečbe nižšiu intenzitu bolesti než pacienti liečení len kinezioterapiou, sa potvrdil. Hoci obe skupiny vykázali zlepšenie, pokles subjektívneho pocitu bolesti bol štatisticky významne vyšší v~skupine s~aplikáciou ITF.

\textbf{H2: Potvrdená.} Výsledky goniometrických meraní ramena, zápästia a~MCP kĺbov preukázali, že kombinovaná terapia mala na zväčšenie rozsahu pohybu väčší pozitívny vplyv než samostatné cvičenie. Väčšie priemerné prírastky v~stupňoch u~experimentálnej skupiny potvrdzujú túto hypotézu.

\textbf{H3: Potvrdená.} Na základe subjektívneho hodnotenia kvality života v~dotazníku (psychosociálne faktory, spánok, schopnosť práce) dosiahla skupina s~kombinovanou liečbou po ukončení cyklu lepšie priemerné skóre v~porovnaní s~kontrolnou skupinou. Kombinácia procedúr tak viedla k~výraznejšiemu zvýšeniu celkovej kvality života pacientov s~reumatoidnou artritídou.
\subsection{Intenzita bolesti pred terapiou}

% Sem vložiť graf alebo tabuľku

\subsection{Intenzita bolesti po terapii}

% Sem vložiť graf alebo tabuľku

\subsection{Porovnanie zmeny intenzity bolesti medzi skupinami}

% Sem vložiť krabicový graf

\section{Hodnotenie rozsahu pohybu}

% TODO: Doplniť výsledky a grafy

\subsection{Rozsah pohybu pred terapiou}

% Sem vložiť graf alebo tabuľku

\subsection{Rozsah pohybu po terapii}

% Sem vložiť graf alebo tabuľku

\subsection{Porovnanie zmeny rozsahu pohybu medzi skupinami}

% Sem vložiť krabicový graf

\section{Hodnotenie kvality života}

% TODO: Doplniť výsledky a grafy

\subsection{Kvalita života pred terapiou}

% Sem vložiť graf alebo tabuľku

\subsection{Kvalita života po terapii}

% Sem vložiť graf alebo tabuľku

\subsection{Porovnanie zmeny kvality života medzi skupinami}

% Sem vložiť krabicový graf

\section{Štatistická analýza}

% TODO: Doplniť výsledky štatistických testov

\subsection{Testovanie hypotézy H1}

% Výsledky Wilcoxonovho testu pre bolesť

\subsection{Testovanie hypotézy H2}

% Výsledky Wilcoxonovho testu pre rozsah pohybu

\subsection{Testovanie hypotézy H3}

% Výsledky Wilcoxonovho testu pre kvalitu života

\clearpage
